\section{Objectifs}
\begin{itemize}
\item Reconnaître un schéma de Bernoulli.
\item Utiliser la loi binomiale.
\item Calculer espérance et variance.
\item Comprendre concentration et loi des grands nombres.
\end{itemize}

\section{1. Épreuve de Bernoulli}
Une épreuve de Bernoulli possède deux issues :
\begin{itemize}
\item succès, de probabilité $p$ ;
\item échec, de probabilité $1-p$.
\end{itemize}

La variable $X$ qui vaut 1 en cas de succès et 0 sinon vérifie
\[
\mathbb E(X)=p,
\]
\[
V(X)=p(1-p).
\]

\section{2. Schéma de Bernoulli}
On répète $n$ épreuves de Bernoulli identiques et indépendantes.

Si $X$ compte le nombre de succès, alors
\[
X\sim\mathcal B(n,p).
\]

Pour $k\in\{0,\ldots,n\}$ :
\[
P(X=k)=\binom nk p^k(1-p)^{n-k}.
\]

\section{3. Espérance et variance}
Pour $X\sim\mathcal B(n,p)$ :
\[
\mathbb E(X)=np,
\]
\[
V(X)=np(1-p).
\]

L'écart-type vaut
\[
\sigma(X)=\sqrt{np(1-p)}.
\]

\section{4. Sommes de variables aléatoires}
Pour des variables admettant une espérance :
\[
\mathbb E(X+Y)=\mathbb E(X)+\mathbb E(Y).
\]

Si $X$ et $Y$ sont indépendantes :
\[
V(X+Y)=V(X)+V(Y).
\]

\section{5. Inégalité de Bienaymé-Tchebychev}
Pour une variable $X$ d'espérance $\mu$ et de variance $V(X)$, pour tout $\delta>0$ :
\[
P(|X-\mu|\ge\delta)\le\frac{V(X)}{\delta^2}.
\]

\section{6. Fréquence et loi des grands nombres}
Pour des variables de Bernoulli indépendantes de même paramètre $p$, la fréquence
\[
F_n=\frac{X_1+\cdots+X_n}{n}
\]
se concentre autour de $p$ lorsque $n$ augmente.

\section{Erreurs fréquentes}
\begin{itemize}
\item Utiliser une loi binomiale sans indépendance ni probabilité de succès constante.
\item Confondre $P(X=k)$ et $P(X\le k)$.
\item Additionner des variances sans indépendance.
\item Interpréter une borne de Tchebychev comme une probabilité exacte.
\end{itemize}

\section{À retenir}
La loi binomiale modélise un nombre de succès dans des essais indépendants ; espérance, variance et concentration décrivent position et dispersion.
